\begin{table}[htbp]
  \caption{Economic Calibration of the Reduced-Form Magnitudes}
  \label{tab:calibration}
  \bigskip
  \centering
  \begin{tabular}{lrrrrrrr}
    \toprule
    City & $s_{WL,d}$ & $11{,}000{\times}s_{WL,d}$ & Households$^{a}$ & $\hat{\pi}s_{WL,d}$ & $\Delta$\%rent$^{b}$ & $\Delta$\$rent$^{c}$ & $\Delta$\%trend$^{d}$ \\
    \midrule
    Kamloops & 0.212 & 2,334 & 292 & 0.174 & 19.0 & \$161 & 13.4 \\
    Vancouver & 0.174 & 1,910 & 239 & 0.142 & 15.3 & \$129 & 10.9 \\
    Prince George & 0.100 & 1,096 & 137 & 0.082 & 8.5 & \$72 & 6.1 \\
    Kelowna & 0.069 & 760 & 95 & 0.057 & 5.8 & \$49 & 4.2 \\
    Quesnel & 0.068 & 743 & 93 & 0.055 & 5.7 & \$48 & 4.1 \\
    \bottomrule
  \end{tabular}
  \par\raggedright\smallskip
  \footnotesize
  Top-five most exposed cities in the main sample. $^{a}$~Potential renter households under the illustrative scenario $(q,h)=(0.25,2)$: $\hat{H}_d = q\times11{,}000\times s_{WL,d}/h$. $^{b}$~$\Delta\%=(\exp(\hat{\pi}s_{WL,d})-1)\times100$ using $\hat{\pi}=0.818$. $^{c}$~$\Delta\$=\Delta\%\times\$848$, the 2015 CMHC average monthly rent used for scale. $^{d}$~Same calculation using the city-trend coefficient $\hat{\pi}=0.594$ from column~(5) of Table~\ref{tab:main_result}. Household counts are scenarios, not estimates of realised post-fire relocations.
\end{table}

